\chapter{Epistemic Conditions for Nash Equilibrium}
\label{sec:p17-epistemic-nash-conditions}

Aumann and Brandenburger ask which information and rationality conditions
make players' conjectures a Nash equilibrium. Their 1995 \emph{Econometrica}
paper distinguishes a conjectured distribution from a conscious randomization
and from an actual chosen action. The two-player result needs mutual
probability-one belief; the many-player result obtains agreement and
independence from a common prior and common probability-one belief in
the full conjectures \parencite[pp.~1161--1164]{AumannBrandenburger1995Epistemic}.
The source's measurable type spaces, complete proofs, numerical examples,
and conceptual qualifications are developed here. A proposed converse in
its note~31 requires the actual-action qualification given below.

\section{The source's interactive belief system}

\begin{hbdefinition}{def:p17-ab-system}{Types, actions, payoffs and probability kernels}
Let \(N=\{1,\ldots,n\}\), with finite \(n\geq2\), and let each \(A_i\)
be a finite nonempty action set. Write \(A=\prod_iA_i\) and
\(A^{-i}=\prod_{j\ne i}A_j\). Give player \(i\) a nonempty measurable
type space \((S_i,\mathcal S_i)\), and let
\((S,\mathcal S)=\prod_i(S_i,\mathcal S_i)\) have the product sigma-field.
A state is a tuple \(s=(s_1,\ldots,s_n)\).
Write \(S^{-i}=\prod_{j\ne i}S_j\) with its product sigma-field.

Each type \(s_i\) specifies an action \(a_i(s_i)\in A_i\), a payoff
array \(g_i(s_i):A\to\mathbb R\), and a probability measure
\(b_i(s_i)\) on \(S^{-i}\). Actions and each coordinate of the
finite payoff array are measurable. Extend the theory \(b_i(s_i)\)
to a probability on \(S\) by fixing the own coordinate:
\[
 p_i(E;s_i)=b_i(s_i)\{t^{-i}:(s_i,t^{-i})\in E\},
 \qquad E\in\mathcal S.
\]
As in the source's general formulation, require
\(s_i\mapsto p_i(E;s_i)\) to be measurable for every \(E\).
All these probabilities are specified at every type, including
types that a later common prior might give probability zero.
The subscript in \(p_i\) makes explicit the player index carried
by the argument \(s_i\) in the source's notation \(p(E;s_i)\).

Write \(a(s)=(a_i(s_i))_i\), \(a^{-i}(s)=(a_j(s_j))_{j\ne i}\),
and \(g(s)=(g_i(s_i))_i\).
For a measurable variable \(X\), the notation \([X=x]\) means its
equality event when that event is measurable. In particular,
\([g=g^*]\) fixes every entry of a specified finite game \(g^*\).
An actual payoff is \(g_i(s_i)(a(s))\), not the whole payoff array.
\end{hbdefinition}

This follows the source's sections~2 and~6
\parencite[pp.~1164--1166, 1173--1174]{AumannBrandenburger1995Epistemic}.
An own type determines an own action and an own payoff function.
Definition~\ref{def:p17-egt-types} instead records choices separately
from types in a fixed game. Here opponents' payoff arrays may vary.
The model describes which choice occurs at a state; it does not prescribe
that choice. Arbitrary measurable spaces are allowed: existence of a
conditional probability kernel is not assumed for every probability law.
The kernels required by this definition are part of the supplied model.

\begin{hbdefinition}{def:p17-ab-conjecture}{Conjectures and rationality}
Player \(i\)'s full conjecture is the probability on all opponents' actions
\[
 \phi^i(s)(c^{-i})=p_i([a^{-i}=c^{-i}];s_i).
\]
Write \(\phi(s)=(\phi^i(s))_i\). Its marginal on \(A_j\), \(j\ne i\),
is \(i\)'s conjecture about \(j\). A full conjecture may correlate
different opponents, and two players may disagree about a third.

Player \(i\) is rational at \(s\) when, for every \(d_i\in A_i\),
\[
 \sum_{c^{-i}}\phi^i(s)(c^{-i})
  \bigl[g_i(s_i)(a_i(s_i),c^{-i})
             -g_i(s_i)(d_i,c^{-i})\bigr]\geq0 .
\]
Let \(R_i\) be this event and \(R=\bigcap_iR_i\).
The own action, payoff array and conjecture all depend only on \(s_i\),
so \(R_i\) depends only on that coordinate. No optimality claim about
actions that a player merely conjectures has yet been imposed.
\end{hbdefinition}

\begin{hbdefinition}{def:p17-ab-knowledge}{Probability-one knowledge and its iterations}
For \(E\in\mathcal S\), put
\[
 K_iE=\{s:p_i(E;s_i)=1\},\qquad K^1E=\bigcap_iK_iE,\qquad
 CK E=\bigcap_{r\geq1}K^rE ,
\]
where \(K^{r+1}E=K^1(K^rE)\) for \(r\geq1\).
Thus \(K^1E\) is mutual knowledge and \(CKE\) common knowledge in the
paper's terminology. They mean probability-one belief, and do not
generally imply that \(E\) is true at the state being assessed.
The definition of \(CKE\) does not silently include \(E\) itself.

A probability \(P\) on \(S\), with marginal \(P_i\), is a
\emph{common prior} when the given kernels satisfy
\[
 P(E\cap\{s:s_i\in B\})
       =\int_B p_i(E;t_i)\,P_i(dt_i)
       \quad(E\in\mathcal S,\ B\in\mathcal S_i,\ i\in N).
 \tag{CP}
\]
This is the precise conditional-distribution identity. It constrains
kernel versions only up to \(P_i\)-null sets; all-type measurability
and the fixed-own-coordinate condition still apply. In finite spaces
it is ordinary conditioning at types with positive prior marginal.
It neither demands positive probability of every state nor supplies
conditioning by division at a null type.
\end{hbdefinition}

\begin{hbtheorem}{thm:p17-ab-measurability}{proved=annex:p17.ab.measurability}{Events and conditional averaging on the full measurable domain}
All events and functions just defined are measurable; the extension
\(p_i(\cdot;s_i)\) is a probability measure. In particular, the
countable iteration defining \(CKE\) is legitimate. Condition (CP) implies
\[
 P(E)=\int p_i(E;t_i)\,P_i(dt_i).
\]
Consequently, for real \(\lambda\), if
\(p_i(E;t_i)=\lambda p_i(F;t_i)\) for every type,
then \(P(E)=\lambda P(F)\). The same conclusion holds when the
equality is only \(P_i\)-almost everywhere.
\end{hbtheorem}

\begin{hbtheorem}{thm:p17-ab-introspection}{proved=annex:p17.ab.introspection}{Own-type certainty and the knowledge identities}
If \(D\) is a measurable event depending only on \(s_i\), then
\(K_iD=D\). Thus \(i\) knows its own action, payoff function, rationality
status and probability assessments. More explicitly,
\[
 K_i[p_i(E;s_i)=\lambda]=[p_i(E;s_i)=\lambda].
\]
Knowledge preserves countable intersections and is monotone.
Furthermore \(CKE\subseteq K_iCKE\) for every \(i\).
For any specified profile \(\varphi=(\varphi^i)_i\) of full conjectures,
\[
 K^1[\phi=\varphi]\subseteq[\phi=\varphi],
 \qquad CK[\phi=\varphi]\subseteq[\phi=\varphi].
\]
These are the source's Lemma~2.6 and Lemmas~4.1, 4.3 and~4.4.
\end{hbtheorem}

\section{Actual actions and two-player conjectures}

\begin{hbtheorem}{thm:p17-ab-preliminary}{proved=annex:p17.ab.preliminary}{The preliminary observation}
At a state \(s\), suppose every player is rational and
\(s\in K^1[a=\alpha]\) for a specified pure profile \(\alpha\).
Then \(a(s)=\alpha\), and \(\alpha\) is a pure Nash equilibrium
of the actual game \(g(s)\). No mutual knowledge of rationality
or of all payoff functions is needed for this implication.
\end{hbtheorem}

\begin{hblemma}{thm:p17-ab-supported-actions}{proved=annex:p17.ab.supported.actions}{The source's positive-conjecture support argument}
Fix a finite game \(g^*\) and a full-conjecture profile \(\varphi\).
Suppose some state belongs to
\[
             K^1[g=g^*]\cap K^1R\cap K^1[\phi=\varphi].
\]
If some other player \(i\ne j\) assigns positive \(\varphi^i\)
probability to player \(j\)'s action \(c_j\), then \(c_j\)
maximizes \(g_j^*(\,\cdot\,,c^{-j})\) averaged under \(\varphi^j\).
This asserts optimality of a positively conjectured action,
not positive conjecture probability of every actually chosen action.
\end{hblemma}

\begin{hbtheorem}{thm:p17-ab-two-player}{proved=annex:p17.ab.two.player}{Aumann--Brandenburger Theorem A}
For two players, if
\[
       K^1[g=g^*]\cap K^1R\cap K^1[\phi=\varphi]\ne\varnothing ,
\]
then \((\varphi^2,\varphi^1)\) is a Nash equilibrium of \(g^*\).
The order reverses because player 2's conjecture is a distribution
on player 1's actions. Neither a common prior nor common knowledge
is assumed. The conclusion concerns these conjectures; it need
not make the actual pure action profile a pure equilibrium.
\end{hbtheorem}

These are the original preliminary observation, Lemma~4.2 and
Theorem~A \parencite[pp.~1167--1168]{AumannBrandenburger1995Epistemic}.
The implication in Theorem~\ref{thm:p17-egt-mutual-conjectures}
uses a fixed game and the beliefs of two selected types. Here the
same support reasoning also tracks payoff uncertainty and the
source's event-based knowledge hypotheses; the models are not
identified simply by renaming their coordinates.

\begin{hblemma}{thm:p17-ab-type-bridge}{proved=annex:p17.ab.type.bridge}{The exact finite bridge to the preceding type model}
If all \(S_i\) are finite and every type uses the same game,
push player \(i\)'s theory forward under
\[
        (s_j)_{j\ne i}\longmapsto(a_j(s_j),s_j)_{j\ne i}.
\]
This produces a model of Definition~\ref{def:p17-egt-types}
with exactly the same action conjectures and rational choices.
For two players the present theorem's mutual-rationality and
mutual-conjecture hypotheses imply the preceding
Theorem~\ref{thm:p17-egt-mutual-conjectures}'s hypotheses.

Conversely, any finite model of that preceding definition embeds,
identifying \(A_i\) with its action set \(C_i\), in a source belief
system with \(S_i=A_i\times T_i\), action the
first coordinate, and theory inherited from the second coordinate.
For two actual types satisfying the preceding theorem's hypotheses,
choose any optimal own action for each. The resulting source state
satisfies the hypotheses of Theorem A. This is a finite fixed-game
bridge, not a universal measurable hierarchy representation.
\end{hblemma}

\begin{hbtheorem}{thm:p17-ab-illustration}{proved=annex:p17.ab.illustration}{The complete source illustration}
For the coordination game
\[
 \begin{array}{c|cc}&c&d\\\hline
 C&(2,2)&(0,0)\\D&(0,0)&(1,1)
 \end{array},
\]
take row types \(C_1,D_1,D_2\) with their lettered actions and
column types \(c_1,d_1,d_2\). Give the following six states prior
mass \(1/6\) each and all other states mass zero:
\[
 (C_1,c_1),\ (C_1,d_1),\ (D_1,c_1),\
 (D_1,d_2),\ (D_2,d_1),\ (D_2,d_2).
\]
Use conditional beliefs from this prior and the displayed game at all types.
At \((D_2,d_2)\), both players are rational and their actions
\((D,d)\) are mutually known. Rationality is not mutually known,
and those actions are not commonly known. Thus the preliminary
observation applies without these stronger assumptions. The annex
calculates every type's expected payoffs.
\end{hbtheorem}

\section{Why a common prior yields agreement and independence}

\begin{hblemma}{thm:p17-ab-local-average}{proved=annex:p17.ab.local.average}{The source's absorbing-event averaging lemma}
Suppose (CP) holds, \(H\subseteq K_iH\), and
\(p_i(E;s_i)=\lambda\) at every state in \(H\). Then
\[
                    P(E\cap H)=\lambda P(H).
\]
The claim includes \(P(H)=0\); no conditional ratio is then taken.
It applies on arbitrary measurable type spaces without assuming
that arbitrary projections of measurable sets are measurable.
\end{hblemma}

\begin{hblemma}{thm:p17-ab-product}{proved=annex:p17.ab.product}{Independence of each coordinate from the rest}
For a probability \(Q\) on the finite product \(A\), suppose
\[
                Q(c)=Q_i(c_i)Q_{-i}(c^{-i})
                     \quad(c\in A,\ i\in N).
\]
Here \(Q_i\) and \(Q_{-i}\) denote the indicated marginal laws.
Then \(Q(c)=\prod_iQ_i(c_i)\). Zeros in the marginals are allowed.
Pairwise independence alone is not the premise.
\end{hblemma}

\begin{hbtheorem}{thm:p17-ab-many-player}{proved=annex:p17.ab.many.player}{Aumann--Brandenburger Theorem B on its exact support domain}
Let \(g^*\) be a finite game and \(\varphi\) a full-conjecture profile.
Suppose a common prior satisfies
\[
 P\bigl(K^1[g=g^*]\cap K^1R\cap CK[\phi=\varphi]\bigr)>0 .
 \tag{AB}
\]
Then all players other than \(j\) agree on a single marginal
\(\sigma_j\) for player \(j\), and
\[
               \varphi^i=\bigotimes_{j\ne i}\sigma_j .
\]
The profile \(\sigma=(\sigma_j)_j\) is a Nash equilibrium of \(g^*\).
This holds for every finite \(n\geq2\); its additional force concerns
three or more players.

More precisely, put \(F=CK[\phi=\varphi]\) and define the conditional
action law \(Q(c)=P([a=c]\cap F)/P(F)\). Then
\[
                   Q=\bigotimes_j\sigma_j .
\]
Thus actual actions lie in the equilibrium supports \(P\)-almost
surely on \(F\). In a finite model this applies to each state of
positive prior mass in \(F\). It does not make a pointwise claim
about arbitrary null states, or identify arbitrary conditional
kernel values there by division.
The agreement and factorization conclusions alone require only (CP)
and \(P(CK[\phi=\varphi])>0\); the game and rationality premises
are used to obtain the Nash conclusion.
\end{hbtheorem}

The two lemmas are Lemmas~4.5 and~4.6, and the proof retains the
original argument through the conditional law on \(F\)
\parencite[pp.~1168--1170]{AumannBrandenburger1995Epistemic}.
The event in (AB) must have positive prior probability; merely
finding a state where its conditions hold is weaker. Common knowledge
concerns each player's \emph{joint} conjecture. Knowledge of individual
marginals does not contain the same information.

\section{All three source tightness examples}

\begin{hbdefinition}{def:p17-ab-marginal-example}{Source Example 5.1: common marginals}
Three players choose \(U,D\), \(L,R\), and \(W,E\). Their payoffs
are given by the two matrices
\[
 \begin{array}{c|cc}
 W&L&R\\\hline U&(1,1,1)&(0,0,0)\\D&(1,0,0)&(1,1,1)
 \end{array}
 \qquad
 \begin{array}{c|cc}
 E&L&R\\\hline U&(0,0,0)&(1,1,1)\\D&(1,1,1)&(0,0,0).
 \end{array}
\]
Each action labels one type with that action and the displayed
own payoff function. The common prior puts mass \(1/4\) on
\((U,L,W),(D,R,W),(U,R,E),(D,L,E)\) and zero elsewhere.
All type marginals are positive; their beliefs are the corresponding
conditional probabilities.
\end{hbdefinition}

\begin{hbtheorem}{thm:p17-ab-marginal-example}{proved=annex:p17.ab.marginal.example}{Agreement of all marginal conjectures is insufficient}
Every type is rational. The game, rationality and all individual
conjectures are commonly known. All individual conjectures are fair.
Their profile is not Nash: against independent fair \(L/R\) and
\(W/E\), row action \(U\) yields \(1/2\) and \(D\) yields \(3/4\).
At \((U,L,W)\), the row player's full conjecture is half \(LW\)
and half \(RE\); the other players do not know this full conjecture.
\end{hbtheorem}

\begin{hbdefinition}{def:p17-ab-chain-game}{Source Examples 5.2 and 5.3: the common game}
The actions are \(H,T\), \(h,t\), and \(W,E\), with payoffs
\[
 \begin{array}{c|cc}
 W&h&t\\\hline H&(1,0,3)&(0,1,0)\\T&(0,1,0)&(1,0,3)
 \end{array}
 \qquad
 \begin{array}{c|cc}
 E&h&t\\\hline H&(1,0,2)&(0,1,2)\\T&(0,1,2)&(1,0,2).
 \end{array}
\]
The first two players play matching pennies; the third receives
3 at matching pairs and 0 at mismatching pairs under \(W\),
and always 2 under \(E\).
\end{hbdefinition}

\begin{hbtheorem}{thm:p17-ab-chain-example}{proved=annex:p17.ab.chain.example}{Source Example 5.2 and its complete finite-depth extension}
Take row types \(H_r,T_r\) and column types \(h_r,t_r\),
\(1\leq r\leq3\), and third-player types \(E_1,W_1,W_2,W_3\).
Their actions are their capital or lower-case letters. Assign
positive prior mass only to the following cells. Each of the four
cells of box \(r\) has the indicated mass, with the third type
specified by its position:
\[
 \begin{array}{c|c|cc}
 r&\text{mass per cell}&\text{diagonal }(H_r,h_r),(T_r,t_r)
                            &\text{off diagonal}\\\hline
 1&9/52&W_1&E_1\\
 2&3/52&W_2&W_1\\
 3&1/52&W_3&W_2 .
 \end{array}
\]
All beliefs are obtained by conditioning, and all own payoff
functions are those of Definition~\ref{def:p17-ab-chain-game}.
Every type is rational. At \((H_2,h_2,W_2)\) the full conjectures
are mutually known and the individual conjectures agree:
fair \(H/T\), fair \(h/t\), and certainty of \(W\). This profile
is not Nash. The unique equilibrium is instead
\((\text{fair},\text{fair},\delta_E)\).

For every positive integer \(m\), a finite extension with
arbitrarily many boxes gives \(m\)-fold mutual knowledge of the
same full conjectures at a positive-prior state, while retaining
agreement, common rationality and the same failure of Nash.
The annex supplies the masses, beliefs and exact distance argument;
no limiting infinite chain is assumed.
\end{hbtheorem}

\begin{hbtheorem}{thm:p17-ab-no-prior}{proved=annex:p17.ab.no.prior}{Source Example 5.3: independence does not follow without the prior}
In the same game use row types \(H_1,T_1\), column types \(h_1,t_1\),
and one third-player type \(W_1\). Every row type assigns half to
each column type and certainty to \(W_1\); every column type assigns
half to each row type and certainty to \(W_1\). Type \(W_1\) assigns
probabilities \(3/8,1/8,1/8,3/8\) to
\((H_1,h_1),(H_1,t_1),(T_1,h_1),(T_1,t_1)\).
All types are rational, and their game, rationality and full
conjectures are commonly known. Individual conjectures agree on
\((\text{fair},\text{fair},\delta_W)\), which is not Nash.
There is no common prior compatible with these beliefs.
\end{hbtheorem}

These are the source's complete Figures~3--7 and Examples~5.1--5.3
\parencite[pp.~1170--1173]{AumannBrandenburger1995Epistemic}.
They distinguish correlation within a conjecture from disagreement
between players, and mutual knowledge from every-order knowledge.

\begin{hbtheorem}{thm:p17-ab-depth-disagreement}{proved=annex:p17.ab.depth.disagreement}{Finite-depth knowledge can also coexist with disagreement}
For every positive integer \(m\), there is a finite three-player
common-prior model, a commonly known zero-payoff game and common
rationality, and a positive-prior state with \(m\)-fold mutual
knowledge of the full conjectures, where players 1 and 2 assign
different probabilities, \(2/3\) and \(1/3\), to player 3 choosing 1.
The explicit construction in the annex supplies the separate
disagreement possibility discussed on source page~1170.
\end{hbtheorem}

\section{Independence, converses and the actual-action qualification}

\begin{hbtheorem}{thm:p17-ab-independent}{proved=annex:p17.ab.independent}{The source's Remark 7.1}
Let \(\sigma\) be a profile of mixed actions. Suppose at some
state the game \(g^*\), rationality, and the following conjecture
facts are mutually known: every other player's conjecture about
player \(j\) equals \(\sigma_j\), and each conjectured opponent
action is independent of the remaining opponents' actions.
Then \(\sigma\) is a Nash equilibrium of \(g^*\).
Neither a common prior nor common knowledge is required here;
agreement and independence have instead been assumed explicitly.
\end{hbtheorem}

\begin{hbtheorem}{thm:p17-ab-converse}{proved=annex:p17.ab.converse}{Realizing an equilibrium and correcting the note-31 converse}
Every Nash equilibrium of a finite game has a finite belief-system
realization satisfying Theorem B, and Theorem A when there are two
players. Each type corresponds to one action in its equilibrium
support, all theories are the appropriate independent opponent laws,
and their product is a common prior. Every supported state is rational
and all full conjectures and payoffs are commonly known.

For a different, pointwise converse, fix \(g^*\), \(\varphi\)
and a mixed profile \(\sigma\) such that
\(\varphi^i=\bigotimes_{j\ne i}\sigma_j\).
Assume at \(s\) that \(g=g^*\) and \(\phi=\varphi\)
are mutually known. For two players the product condition is
automatic with \(\sigma=(\varphi^2,\varphi^1)\).
Then rationality is mutually known at \(s\) if and only if
\(\sigma\) is Nash in \(g^*\) \emph{and} each actual \(a_i(s_i)\)
is a best response to \(\sigma_{-i}\).
Membership of every actual action in its positive equilibrium support
is a sufficient, stronger replacement for the last condition.
For the many-player common-prior setting, the product condition
follows already from \(P(CK[\phi=\varphi])>0\), as above.
For the preliminary observation, if the specified actions \(\alpha\)
and game \(g^*\) are mutually known, actual rationality of all
players is equivalent to \(\alpha\) being Nash in \(g^*\).
Mutual rationality need not follow, as the source illustration shows.

Nash conjectures alone do not imply mutual rationality under these
premises. This corrects the converse asserted in the original
note~31 without changing the valid Theorems A or B.
\end{hbtheorem}

\begin{hbtheorem}{thm:p17-ab-support-failures}{proved=annex:p17.ab.support.failures}{Explicit failures of the removed hypotheses}
Finite models in the annex establish each of the following:
Nash conjectures and commonly known payoffs/conjectures can coexist
with irrational actual actions; actual rationality cannot replace
mutual rationality; knowledge of one's own payoff function cannot
replace mutual knowledge of the game; and nonemptiness cannot replace
the positive-prior premise in (AB).
The examples specify every type, action, payoff, belief and prior
where one exists. Null states remain part of the model and are
identified explicitly.
\end{hbtheorem}

\begin{hbtheorem}{thm:p17-ab-knowledge-of-equilibrium}{proved=annex:p17.ab.knowledge.of.equilibrium}{Knowing a specified profile is an equilibrium}
At a state meeting Theorem A's conditions, each player knows that
the specified conjecture profile is Nash in the game being played.
The same holds at states in the event of Theorem B's hypothesis.
Here the event is that the \emph{fixed} resulting profile is Nash
in \(g(s)\), not that a newly evaluated conjecture at every possible
state has been identified. Higher-order knowledge need not follow.
The preliminary observation alone need not even give first-order
mutual knowledge of this equilibrium fact; adding mutual knowledge
of the payoff functions does give it.
\end{hbtheorem}

\section{Interpretation and use}

The source's belief system is a descriptive mathematical language:
it can state what actions and beliefs coexist and whether its optimality
criterion holds. It does not explain why someone has those beliefs,
recommend an action, or assert that real players satisfy the conditions.
Including actions in the state does not remove freedom of choice;
choosing differently gives a different state description. This differs
from the one-person decision use of a state as something unaffected by
the decision maker. Harsanyi's related game framework uses types for
payoffs without requiring that a type also determine an action
\parencite[pp.~1174--1175]{AumannBrandenburger1995Epistemic}.

Probability-one belief is weaker than impossibility of error. The source
uses the former, so replacing its knowledge premises by a genuine
certainty condition that implies them preserves its sufficient
implications. It does not make the event operators factive. Own-type
certainty includes certainty of own action and own payoff function.
For utility affected by unknown external circumstances, the source
suggests interpreting its array as expected payoff at each action
profile; this is an interpretation, not a theorem constructing such
an array from every uncertain-utility model. If conscious randomization
is modeled, knowledge may concern the chosen mixture while leaving
its realization unknown. That is a different action domain; the present
formal results retain finite actions interpreted as definite choices
\parencite[pp.~1175--1176]{AumannBrandenburger1995Epistemic}.

Knowledge of conjectures is substantive. The source's car-purchase
story suggests that a customer might know the statistical information
a seller uses to form a conjecture about customers, even without the
seller knowing that the customer knows it. The story illustrates a
possible context for mutual rather than common knowledge; it supplies
no measured frequency or claim that such information is usually available.
Likewise, describing a belief system in one common mathematical
language is different from assuming common knowledge of the actual
game or payoff functions. The paper revises earlier informal claims
that a real, externally specified game model must always be commonly
known \parencite[pp.~1176--1177]{AumannBrandenburger1995Epistemic}.

Theorem A weakens the common-knowledge premises in the two-player
results discussed by Tan--Werlang and Brandenburger--Dekel. Those
presentations use their own definitions and additional claims.
The paper also distinguishes its local question about a particular
state's conjectures from Aumann's global common-prior characterization
of a distribution of action profiles as a correlated equilibrium.
Common rationality in an entire model does not make that distribution
a product. Theorem~\ref{thm:p17-ab-independent} instead assumes the
agreement and independence it needs. The source regards that route as
less informative about the epistemic origin of independence; its
note~32 distinguishes it from the earlier Brandenburger--Dekel
conditions, which include a different prior-consistency assumption
\parencite[pp.~1177--1179]{AumannBrandenburger1995Epistemic}.

\begin{hbmethod}{def:p17-ab-computation}{Checking a finite epistemic model}
Specify every type's action, payoff array and joint opponent-type
probability. Marginalize to full action conjectures and compute all
pure-action expected payoffs. Mark \(R_i\), the specified game event,
and the specified full-conjecture event. For an event, mark a type
as knowing it exactly when every state receiving positive probability
at that type belongs to the event. Iterate the mutual operator and
intersect all attained levels to obtain common knowledge; finite
set iteration can equivalently use all reachable states.

If using Theorem B, provide a common prior and verify its conditional
identities at every positive-marginal type. Sum the prior over the
actual hypothesis event before dividing by \(P(F)\). Compute \(Q\),
its marginals and the Nash support inequalities. Return conjectures,
actual actions and any null-state qualifications separately.
This checks a supplied model; it does not infer its epistemic
assumptions from observed choices.
\end{hbmethod}

The finite reachability computation is justified in
Theorem~\ref{thm:p17-ab-introspection}'s proof. The results retain
arbitrary measurable type spaces because only finite action sums,
countable event operations and the explicitly supplied conditional
identities are used. A universal space representing every possible
belief hierarchy, a one-player reformulation of another observer's
knowledge, and general learning or equilibrium-selection processes are
separate constructions. For one player there are no opponent
conjectures from which to recover an own mixed action; ordinary
maximization gives the applicable equilibrium condition directly.
